% r06_fill_HS: the heavy-shell estimates (HS), Theorem rot:hyp:HS: the two shell theorems and the proof of (HS) from them.
% Input from sec/r06_rot_shell2.tex in place of the proof of \Cref{rot:hyp:HS}.
\providecommand{\Ish}{\ensuremath{I^{\sharp}}}
\providecommand{\Ifull}{\ensuremath{I}}
\providecommand{\Scap}{\ensuremath{\tilde S^{\mathrm{cap}}}}
\providecommand{\chicap}{\ensuremath{\tilde\chi^{\mathrm{cap}}}}
\providecommand{\kz}{\ensuremath{\varkappa_0}}

\subsubsection*{The shell model and the stopping-scale integrals}

Fix the exterior-shell configuration of \Cref{rot:hyp:HS}: the torus $T=D_{J'}$, the box $\Hboxp=\Qb{b'}$, both aligned with the blocks of the final scale as \Cref{rot:hyp:HS} requires (block alignment and grid offset), the plateau height $n\ge n_*=1$ and the shell $U=T\setminus\Hboxp$. The \emph{shell model} at shell mass $\varkappa\in[0,\infty)$ is King's walled torus model of \Cref{rot:def:tori}(c) with the wall $n\psi+\varkappa\mathbf 1_U$: its matter operator is $-\Delta_A+m^2+\dm_\eps+2n\psi+2\varkappa\mathbf 1_U$, and $\varkappa\mathbf 1_U$ is kept inside the matter covariance of every block-spin step. King's block-spin flow \cite[\S3]{K1} is run on $\Lk$, $\eps=\eps_k=5^{-k}$, through $r$ steps down to the \emph{stopping scale} $\eta_0\defeq5^{r-k}\in5^{-\N}$, at which the plateau is massive in block units (no quantitative heaviness of the plateau is used; only $n\ge1$ enters). The shell mass in block units at the stopping scale is $w_r\defeq\varkappa\eta_0^2$; the \emph{frozen regime} is $w_r\ge c_0$ for a fixed $c_0\in(0,c_0^*]$, and the threshold of \Cref{rot:hyp:HS} is
\[
\kz=\varkappa_0(\eta_0)\defeq\hat c_0/\eta_0^2,\qquad\hat c_0\defeq\max(c_0,1).
\]
The output of the flow is the effective action $S^{(r),\varkappa}(A_r,\phi_r;zf,0)$ on the block fields of the coarse lattice $5^r\eps\,\Z^3/P_{J'}$, in King's unit scale, with the curvature source $zf$ carried as in \cite[(3.20)]{K1}, and the small-field indicator $\chi_r$ of \cite[(3.2)]{K1}; the shell model carries in addition two tests at the stopping scale, a \emph{near-shell test} $\tau_r$ on the block fields of $U$ and a \emph{first-layer test} $\lambda_r$ on the block fields of the first light layer of $\Hboxp$, both with thresholds that depend on $w_r$ through increasing functions, and a weakened gradient family $\chi_r^\flat\ge\chi_r$ with no clause on the bonds inside $U$ or crossing $\partial U$. The stopping-scale integrals are
\[
\Ish_\varkappa(z)\defeq\int(dA_r)(d\phi_r)\,\chi_r\tau_r\lambda_r\,e^{-S^{(r),\varkappa}},\qquad
\Ifull_\varkappa(z)\defeq\int(dA_r)(d\phi_r)\,\chi_r^\flat\tau_r\lambda_r\,e^{-S^{(r),\varkappa}}\ \ge\ \Ish_\varkappa(z).
\]
For $\varkappa\ge\kz$ the thresholds of $\tau_r$, $\lambda_r$ are dominated by the $\varkappa$-free caps $C_7p_r$ on $U$ and $C_8p_r$ on the first layer ($p_r=p(\eta_0)$ King's small-field function), and the \emph{capped} integral $\Icap_\varkappa(z)$ of \Cref{rot:hyp:HS} is $\Ifull_\varkappa$ with $\tau_r$, $\lambda_r$ replaced by the indicators of these caps. Under the standing smallness $\eta_0\le5/(\lambda C_7^4)$ the cap on $U$ implies the block-sup clause of $\chi_r^\flat$ there, so the capped test set is a product: with $b\defeq(A_r,\phi_r|_{\Hboxp})$ the \emph{exterior block data} and $\phi_U\defeq\phi_r|_U$,
\[
\chicap\defeq\{b:\ \chi_r^\flat\lambda^{\rm cap}_r(b)=1\},\qquad
\Scap_\varkappa(b;z)\defeq-\ln\int\prod_{y\in U^{(r)}}d\phi_r(y)\ \mathbf 1\{|\phi_r|\le C_7p_r\text{ on }U^{(r)}\}\,e^{-S^{(r),\varkappa}(A_r,\phi_r;zf,0)},
\]
so that $\Icap_\varkappa(z)=\int_{\chicap}e^{-\Scap_\varkappa(b;z)}\,db$ exactly: this is the measure space $(\tilde\chi,db)$ and the reduced small-field action of \Cref{rot:hyp:HS}, and the test set $\chicap$ is one and the same for every $\varkappa\ge\kz$, which is the threshold clause of that theorem. Every statement below is uniform in $\eps\le\eps_0(\eta_0)$, in $\varkappa$ on the range stated, in $z\in[-1,1]$, and holds for the model with the free-boundary counterterm $\dm_\eps$ of \Cref{def:s02:acube}: the torus-convention counterterm of \cite{K1} differs from it by a number $\Delta_T(\eps)$ bounded uniformly in $\eps$ (\Cref{rot:prop:CT}(b) above), and every bound of the shell analysis holds with constants depending on $\sup_\eps|\Delta_T(\eps)|$ as well.

\begin{proposition}[the two-sided bound with the shell inside]\label{rot:prop:SW1}
Let $L=5$, the couplings $e,\lambda,m,\lmu$ in King's range, $(\Hboxp,T)$ block-aligned with the grid offset as in \Cref{rot:hyp:HS}, $n\ge1$, and $f$ as in \Cref{rot:hyp:HS}. There are $\eta_1>0$, $\sigma>0$, $c>0$ and, for every $\eta_0\in5^{-\N}$ with $\eta_0\le\eta_1$, a spacing bound $\eps_0(\eta_0)>0$ and a constant $C<\infty$ depending on $(\lambda,e,m,\lmu,|T|,n,f,\Hboxp,c_0)$ but not on $\eta_0$, $\varkappa$ or $\eps$, such that for all $\eps\le\eps_0(\eta_0)$ and $z\in[-1,1]$, with
\[
\sfa(\eta_0)\defeq C\eta_0^{\sigma}|T|,\qquad\sfb(\eta_0)\defeq e^{-c\,p(\eta_0)^2+C|T|},
\]
\begin{align*}
Z_\varkappa(zf,0)&\ \ge\ e^{-\sfa}\,\Ish_\varkappa(z)&&(\varkappa\in[0,\infty)),\\
Z_\varkappa(zf,0)&\ \ge\ e^{-\sfa}\,\Icap_\varkappa(z)&&(\varkappa\in[\kz,\infty)),\\
Z_\varkappa(zf,0)&\ \le\ e^{\sfa}\,\Ifull_\varkappa(z)+\sfb\,Z_\varkappa(0,0)\ \le\ e^{\sfa}\,\Icap_\varkappa(z)+\sfb\,Z_\varkappa(0,0)&&(\varkappa\in[0,\infty)).
\end{align*}
In particular $\sfa,\sfb\to0$ as $\eta_0\to0$ at fixed $(n,|T|)$, uniformly in $\varkappa$ and $\eps\le\eps_0(\eta_0)$, and clause \textup{(i)} of \Cref{rot:hyp:HS} holds for $\varkappa\ge\kz$ with the capped integral.
\end{proposition}

\begin{proof}
The two lines are King's Theorem~3.1 \cite[(3.3)--(3.4)]{K1}, that is Ba{\l}aban's lower bound \cite[\S\S3--4]{Ba1} and upper bound \cite[\S\S2--3]{Ba2} with the renormalisation of \cite{Ba3}, re-run for the shell model, at block factor $5$ (\Cref{thm:a06:BALA}) and with the remainder of the upper line measured against $Z_\varkappa(0,0)$. The proof has four layers, each uniform in $\varkappa\in[0,\infty]$.

\emph{(1) Exact algebra.} The block-spin transformation is a Gaussian identity valid for every positive matter covariance: with $\varkappa\mathbf 1_U$ a non-negative diagonal operator, every object of \cite[(2.13)--(2.17), (3.14)]{K1} is defined, the composition law of the steps, the scale decomposition of the fluctuation covariance into slices, the minimiser identity for the background $\phi^{(j)}_\varkappa=a_j(M^\varkappa_j)^{-1}Q_j(A)^*\phi_j$ and the Schur-complement monotonicity in $\varkappa$ hold verbatim, and every scalar line of the flow is a leg of block data (the slice covariances are $\ell\,C^{(j),\varkappa}\ell^\dagger$ with legs $\ell$). A two-line positivity argument inserts the near-shell test at every scale at no cost.

\emph{(2) Operator bounds.} King's Theorem~3.3 \cite[(3.5)--(3.8)]{K1} for the shell model in the forms consumed: exponential decay, sup and first-difference bounds for the step operators $G_j$, $\Delta^{(j)}$, $C^{(j)}$ everywhere, uniformly in $\varkappa$ (coercivity, since $\varkappa\mathbf 1_U\ge0$ adds to the quadratic form; $\ell^2$ decay by conjugation, which $\varkappa\mathbf 1_U$ commutes with; Loewner monotonicity for the unit-lattice kernels; Kato's inequality with the free Laplacian for the fine sup bounds); the derivative clause for the background of small-field data; H\"older bounds on the light side; and on the heavy side the amplitude bounds of a barrier argument (a constant supersolution of mass $\sqrt{w_j}$ in block units), with the first-difference bounds at the heavy points paired with the amplitude of the other leg of the pair. The pointwise difference bounds of the line kernels, $|D_AG(x,y)|\le C|x-y|_\eta^{-2}e^{-c|x-y|}$ and, with one difference at each end of one line, $|D^x_AD^y_AG(x,y)|\le C|x-y|_\eta^{-3}e^{-c|x-y|}$, uniformly in $\varkappa$ on both sides of $\partial U$, follow from the scale decomposition of (1), one difference per variable carried by its own leg, the leg bounds and a geometric sum over the scales. On the Neumann sub-domains of Ba{\l}aban's splitting the first-difference and H\"older clauses are used at points at distance at least Ba{\l}aban's localisation radius from the complement, where they hold; the convexity of $\partial U=\partial\Hboxp$ enters here, since at a re-entrant edge the first-difference bound fails uniformly in the spacing.

\emph{(3) Perturbative bounds.} Ba{\l}aban's renormalisation bound \cite[Prop.~1, (1.33)]{Ba3} for the renormalised families of \cite[\S3.5]{K1} of order at most $13$ with at least one external field, with the shell model's propagators and legs, in Ba{\l}aban's product form with a $\varkappa$-adapted norm on the scalar external fields, and the restricted form of King's bound \cite[(3.38)]{K1} on the field-dependent families; the Taylor step at a divergent two-scalar-leg subgraph acts on the photon legs, whose norms do not see the shell mass, with the remainder controlled by the current two-point kernel bound $|\Pi^\varkappa_j(x,y)|\le C|x-y|^{-4}$ from the difference bounds of (2); the two-photon-leg and mixed classes are combined by the Ward--Takahashi identities, which hold for the shell model since $\varkappa\mathbf 1_U$ is gauge invariant, used for full step-$j$ families and never for a fixed ordering of the lines; a mismatch vertex $\mathsf M$ (the difference of the shell-adapted and the bulk subtraction inside a renormalised pair) is a two-leg vertex of bounded coefficient. The block-data form of the bound, with King's legs $\phi^{(r)}_\varkappa$ and the small-field data $|\phi_r|\le p$, $|\partial A_r|\le p$, $|D_{A_r}\phi_r|\le p$ of \cite[(3.2)]{K1}, is derived from the product form on the photon legs by a local gauge transformation that removes the constant part of the background before the field itself is bounded, King's data satisfying $|A_r|\le p\,(\lmu\ell_{r-1})^{-1}$ but no uniform bound; the same reduction gives the light-side first moment of the photon two-leg kernel, summed against the background and its first differences, which the near-face analysis of \cite[\S5.2]{K2} transplanted to the shell model consumes. \moving{proof of the bound, uniform in the shell mass, on the first-order term in the background field of the determinant of the step fluctuation Gaussian on the Neumann sub-domains of Ba{\l}aban's splitting, where the large-field analysis of the upper bound uses it: being completed}

\emph{(4) The two bounds.} The lower bound is \cite[\S\S3--4]{Ba1}: restriction to the small-field region by Jensen's inequality, the inductive hypothesis on King's Appendix class \cite[(A.1)--(A.8)]{K1} for the shell model (with a background-gradient and a block-mismatch clause, which compose from step to step), the cumulant bound for the step Gaussian with $\varkappa$-free constants (\Cref{rot:thm:R1}), and the final step \cite[(3.66)--(3.69)]{Ba1}; the restriction to the capped set at the stopping scale is a positivity step inside the class, which gives the second line for $\varkappa\ge\kz$. The upper bound is \cite[\S\S2--3]{Ba2}: the large/small-field split at every scale, with the near-shell test paid at every scale and the first-layer test at the stopping scale, the sum over large-field configurations $\prod_j\sum_\Lambda e^{-p(\ell_j)|\Lambda|}\le e^{C|T|}$, the propagator-domain comparisons of \cite[Prop.~2.6]{Ba2} in the form \Cref{hyp:k01:BALint} with the regrouping \Cref{hyp:k01:RENorb}, and the continuation of the large-field terms past the stopping scale, which under block alignment runs the shell model at the scales below $\eta_0$; the remainder is $\sfb\,Z_\varkappa(0,0)$ in normalisation-free form by the ultraviolet-stability gap, as in \eqref{eq:k01:H0rel}. The upper capped line follows from $\Ifull_\varkappa\le\Icap_\varkappa$ (the cap tests dominate $\tau_r$, $\lambda_r$ pointwise as indicators and the integrand is positive); the lower capped line is the positivity step inside the class just described, the capped test set lying in King's Appendix class for $\varkappa\ge\kz$.
\end{proof}

\begin{proposition}[the $\varkappa$-Cauchy property of the reduced action]\label{rot:prop:SW3}
In the setting of \Cref{rot:prop:SW1} fix $\alpha_0\in(0,\frac12]$. There is $\eta_2\le\eta_1$ such that for every $\eta_0\in5^{-\N}$ with $\eta_0\le\eta_2$ there is $C_{\mathfrak c}(\eta_0)<\infty$, depending on $(T,\Hboxp,n,\lambda,e,m,\lmu,\alpha_0,c_0,\eta_0)$ and on $f$ but not on $\eps$, $\varkappa$, $z$, such that for all $\eps\le\eps_0(\eta_0)$, all $\kz\le\varkappa\le\varkappa'<\infty$ and all $z\in[-1,1]$,
\[
\sup_{b,b'\in\chicap}\big|\Delta_z(b)-\Delta_0(b')\big|\ \le\ 2R(\varkappa,\varkappa')\ \le\ 2\,r(\varkappa;\eta_0),\qquad
R(\varkappa,\varkappa')\defeq\int_{\varkappa\eta_0^2}^{\varkappa'\eta_0^2}\theta(w)\,dw+\tau(\varkappa),
\]
\[
\theta(w)\defeq C_{\mathfrak c}\,w^{-1}w^{-(1-\alpha_0)/2}(1+\log w)^2,\qquad
\tau(\varkappa)\defeq C_{\mathfrak c}\,(\varkappa\eta_0^2)^{-\frac12+\frac{\alpha_0}4},\qquad
r(K;\eta_0)\defeq C'_{\mathfrak c}\,(K\eta_0^2)^{-(1-\alpha_0)/2}\big(1+\log(K\eta_0^2)\big)^2,
\]
with $\Delta_z(b)=\Scap_{\varkappa'}(b;z)-\Scap_\varkappa(b;z)$ as in \Cref{rot:hyp:HS}; $r(K;\eta_0)\to0$ as $K\to\infty$ at fixed $\eta_0$. This is clause \textup{(ii)} of \Cref{rot:hyp:HS}.
\end{proposition}

\begin{proof}
The statement is derived from a bound on King's full explicit action in all block fields: there is a field-free function $\varkappa\mapsto c_\varkappa$ on $[\kz,\infty)$, independent of $z$, with
\[
\sup_{(A_r,\phi_r)\in\supp(\chi^\flat_r\tau^{\rm cap}_r\lambda^{\rm cap}_r)}\big|S^{(r),\varkappa}-S^{(r),\varkappa'}-(c_\varkappa-c_{\varkappa'})\big|\ \le\ R(\varkappa,\varkappa'),
\]
and on the capped set, whose $U$-section is the same cube $\{|\phi_U|\le C_7p_r\}$ for every exterior datum $b$ and every $\varkappa$, the reduced actions inherit the bound with the same $c_\varkappa$ (the logarithm of an integral of $e^{-S}$ over a fixed set moves by at most the sup of the change of $S$); taking the mixed difference removes $c_\varkappa$. The full-field bound is proved in derivative form: $\partial_\varkappa S^{(r),\varkappa}$ is King's sum of renormalised graphs with one insertion of $\partial_\varkappa$ of the matter operator, that is of $2\mathbf 1_U$ on a line, plus the $\varkappa$-derivative of the block-free part of the background determinant. For an insertion on a line with a scalar connection to the rest of the graph the heavy matter on $U$ pays $w_j^{-1/2}$ per deep end at every frozen scale (the barrier bounds of layer (2) of the proof of \Cref{rot:prop:SW1}, with the one-loop $z^2$-term and its positions in $U$ treated in both regimes); for an insertion whose matter component through it is closed and connects only through photon lines, which do not see the shell mass, the members with the same full lines are grouped into one class, one gauge identity is applied to the class, and only then is it expanded into King's members with the background rebased for each, the base layer being a first-moment bound, uniform in $\varkappa$, for the two-photon kernel built from full lines. In the units of the stopping scale, with $\hat\varkappa=\varkappa\eta_0^2$, the sum over the members less its field-free part is at most $\theta(\hat\varkappa)$ per unit of $\hat\varkappa$, integrable at infinity; the H\"older remainders of the Taylor step for lines used as legs and the lattice regime of the scale sum give the exponent $(1-\alpha_0)/2$. The block-free part of the determinant has the monotone tail $0\le\Phi^{\rm bf}(\varkappa)-\Phi^{\rm bf}(\infty)\le\tau(\varkappa)$, with the bracket for the loops winding through $U$ (whose contribution is positive even at zero background curvature) bounded, under block alignment, by $C\hat\varkappa^{-1/2}$. Integrating in $\varkappa$ gives $R(\varkappa,\varkappa')$, and the sum over the dyadic intervals of $\hat\varkappa$ gives $r(K;\eta_0)$. The perturbative inputs are those of layer (3) of the proof of \Cref{rot:prop:SW1}, the restricted form of \cite[(3.38)]{K1} being re-run for the shell vertex set; the block-data form is consumed at the step marked there.
\end{proof}

\begin{proof}[Proof of \Cref{rot:hyp:HS}]
Take $\eta_1$ of the theorem to be the $\eta_2$ of \Cref{rot:prop:SW3}, $\sfa,\sfb$ and $\eps_0(\eta_0)$ as in \Cref{rot:prop:SW1}, $\varkappa_0=\kz$, $r(\cdot;\eta_0)$ as in \Cref{rot:prop:SW3}, and $(\tilde\chi,db)=(\chicap,db)$ with $\tilde S_\varkappa=\Scap_\varkappa$. Clause (i) is the second and third lines of \Cref{rot:prop:SW1} for $\varkappa\ge\kz$; clause (ii) is \Cref{rot:prop:SW3}; the test set does not depend on $\varkappa$ on $[\kz,\infty)$, as stated.
\end{proof}
